\chapter{Project State Summary - May 5, 2026}

\section{Objective}
The primary goals addressed recently were the implementation of \textbf{Auto K-Means and DBSCAN}, resolving simulation stability issues, performing a general codebase cleanup, and standardizing the \textbf{Birth-Death simulation modules} for consistent scientific modeling.

\section{Completed Work}

\subsection{1. Dead Agent Archive \& NetCDF Storage (April 21st)}
\begin{itemize}
\item \textbf{Objective}: Implement a performance-conscious system for archiving deceased agents for historical analysis.
\item \textbf{Work Done}:
\begin{itemize}
\item \textbf{Fortran Backend}: Added \texttt{store\_dead\_agents} flag to \texttt{world\_container} and logic to archive agent data before deletion.
\item \textbf{Asynchronous Pipeline}: Built an async Python writer using \texttt{multiprocessing.Queue} to persist data to NetCDF without blocking the simulation thread.
\item \textbf{UI Integration}: Added a checkbox in the "Full Simulation" tab and status indicators for the writer queue.
\item \textbf{Bug Fixes}: Resolved segfaults in the F2PY interface by refactoring agent extraction into a clean dataclass-based pipeline.
\end{itemize}
\end{itemize}

\subsection{2. New Module Pattern \& World Data Access (April 23rd)}
\begin{itemize}
\item \textbf{Objective}: Standardize how new agent-logic modules are added and how they interact with global state.
\item \textbf{Work Done}:
\begin{itemize}
\item \textbf{Accumulator Reset}: Established a mechanism for \texttt{t\_tick\_accumulators} that automatically reset at the end of each simulation tick.
\item \textbf{Module Registration}: Defined a clear 3-step pattern for adding modules (Fortran ID -> Python Interface -> SpawnPointEditor registration) following the \texttt{reviewed\_agent\_motion} template.
\item \textbf{World Access}: Identified and documented accessible agent variables via the \texttt{world} object to facilitate custom agent-logic development.
\end{itemize}
\end{itemize}

\subsection{3. DBSCAN, Convex Hull \& Performance Metrics (April 28th)}
\begin{itemize}
\item \textbf{Objective}: Expand clustering options, improve cluster integrity, and track simulation efficiency.
\item \textbf{Work Done}:
\begin{itemize}
\item \textbf{DBSCAN Implementation}: Added point-based and grid-based DBSCAN clustering in Fortran.
\item \textbf{UI Controls}: Added dynamic control over DBSCAN parameters (\texttt{eps}, \texttt{minPts}) directly from the Python GUI.
\item \textbf{Convex Hull Fill}: Integrated a convex hull post-processing step into the grid-based voting mechanism to eliminate holes and noise gaps within identified clusters.
\item \textbf{Performance Monitoring}: Implemented \texttt{avg\_ms\_per\_tick} tracking. The average time per tick is now displayed in the window title and is available as a plottable variable in the time-series UI.
\end{itemize}
\end{itemize}

\subsection{4. Unified Density Smoothing \& Global Metrics (April 29th)}
\begin{itemize}
\item \textbf{Objective}: Standardize density smoothing across all clustering modes and enable real-time plotting of global states.
\item \textbf{Work Done}:
\begin{itemize}
\item \textbf{Unified Smoothing}: Replaced redundant local smoothing in Auto K-Means with a global \texttt{human\_density\_smoothing\_radius} parameter.
\item \textbf{Iterative Smoothing}: Added \texttt{human\_density\_smoothing\_iterations} to allow for aggressive, multi-pass global smoothing.
\item \textbf{Global Metrics}: Exported \texttt{t\_dynamic\_state} and \texttt{t\_tick\_accumulators} to Python. Added variables like \texttt{k\_fertility (sim)}, \texttt{phi\_death\_acc (sim)}, and various \texttt{death\_* (sim)} debug counters to the UI plotting engine.
\item \textbf{Dual-Axis Fix}: Resolved a PyQtGraph closure bug where secondary axes caused the plot canvas to collapse.
\end{itemize}
\end{itemize}

\subsection{5. Birth-Death Refactoring \& Logic Consolidation (May 4th)}
\begin{itemize}
\item \textbf{Objective}: Standardize parameter naming and consolidate biological logic for easier scientific review.
\item \textbf{Work Done}:
\begin{itemize}
\item \textbf{Parameter Schema Migration}: Replaced legacy \texttt{b1-b4} identifiers with standardized ecological parameters: \textbf{\texttt{r}} (growth rate), \textbf{\texttt{NC}} (carrying capacity), \textbf{\texttt{Kmin}}, and \textbf{\texttt{Kmax}}.
\item \textbf{Module Consolidation}: Migrated birth/death logic from the obsolete \texttt{mod\_birth\_death\_new.f95} into the reviewed and consolidated \texttt{mod\_reviewed\_modules.f95}.
\item \textbf{Build System Sync}: Updated the dependency graph and Python interface to support the new parameter schema and modular structure.
\end{itemize}
\end{itemize}

\subsection{6. Architectural Refinement \& Critical Fixes (April - May)}
\begin{itemize}
\item \textbf{K-Means Farthest-First Traversal}: Replaced the peak-value search in K-Means with a \textbf{Farthest First Traversal (K-means++)} subroutine to ensure geographically remote populations are isolated uniformly.
\item \textbf{Cluster ID 1-Indexing}: Fixed a visualization bug where Cluster ID 0 was being masked as noise; the system now uses 1-based indexing for all valid clusters.
\item \textbf{Rendering \& Matrix Edge Geometry}: Resolved line rendering issues by translating cluster maps directly through NumPy boolean slicing filters and "burning" 1-pixel white boundaries into the RGBA output matrices.
\item \textbf{3D Visualization}: Replaced unstable OpenGL 3D views with Matplotlib 3D surfaces in the \texttt{compare\_smoothing.py} tool for robust cluster verification in headless environments.
\item \textbf{F2PY Namespace Hotfix}: Fixed a critical bug where localized imports of \texttt{mod\_python\_interface} shadowed the extension module, causing \texttt{AttributeError} crashes in the update loop.
\end{itemize}

\section{TODO later}
\begin{itemize}
\item \textbf{[LATEX/DOCS]}: Add an entry explaining the Farthest-First Traversal algorithm.
\item \textbf{[PYTHON/DOCS]}: Document that visual rendering loops for 2D and 3D are fundamentally separate.
\item \textbf{[PYTHON/DOCS]}: Mark that the backend expects exactly 4-channel RGBA arrays for map updates.
\end{itemize}

\subsection{7. Cluster Population Control \& Multi-Population Graphing (May 13th)}
\begin{itemize}
\item \textbf{Objective}: Resolve population convergence discrepancies between global and cluster-based modules and improve the Python graphing capabilities.
\item \textbf{Work Done}:
\begin{itemize}
\item \textbf{Divergence Analysis}: Identified that the cluster-based module (\texttt{new\_birth}) was severely under-suppressing births compared to the global (\texttt{reviewed\_birth}) module. This was caused by the cluster module relying on the per-cluster accumulator \texttt{n\_alive\_acc(jp)}, which structurally misses agents not assigned to any cluster (e.g. \texttt{c\_idx = 0}), resulting in massive population under-estimation.
\item \textbf{Debug Tooling}: Implemented a global-vs-cluster diagnostic summation script within \texttt{python\_interface.f95} to track the exact magnitude of unassigned agents escaping the control logic.
\item \textbf{Dynamic UI Plotting via Population Filters}:
\item Enhanced \texttt{mod\_python\_interface.get\_dynamic\_state\_stats} and \texttt{get\_cluster\_info} to optionally accept and return stats based on a specific \texttt{jp} population index.
\item Added a dedicated \texttt{pop:} spinbox to the Python GUI (in \texttt{application.py}) that feeds into the graph configuration.
\item Implemented \textbf{"Dominant Population" (\texttt{pop = -1}) logic}, allowing graphs to dynamically query the backend tick-by-tick and automatically select and plot statistics (\texttt{NC}, \texttt{n\_agents}, \texttt{k\_fertility}) specifically for whichever population currently has the most agents within the targeted cluster or globally.
\end{itemize}
\end{itemize}

\subsection{8. Performance Metrics HUD \& Population-Specific Mating (May 26th)}
\begin{itemize}
\item \textbf{Objective}: Add real-time wall-clock profiling to identify simulation bottlenecks, and introduce optional reproductive isolation between populations.
\item \textbf{Work Done}:
\begin{itemize}
\item \textbf{Fortran Instrumentation}: Wrapped each phase of \texttt{step\_simulation} in \texttt{python\_interface.f95} with \texttt{system\_clock} calls to measure wall-clock time for permanent modules, active modules, compaction, grid/density, and clustering.
\item \textbf{API}: Exposed \texttt{set\_performance\_timing\_enabled} and \texttt{get\_performance\_stats} to Python via \texttt{mod\_python\_interface}. Implemented as optional with \textbf{zero overhead} when disabled.
\item \textbf{GUI HUD}: Added a ``Show Performance Metrics'' checkbox in the ``View Editor'' tab and a glassmorphic tech-cyan overlay in the simulation window displaying per-phase timing breakdowns.
\item \textbf{Population-Specific Mating (\texttt{allow\_across\_populations})}: Added a new config flag \texttt{allow\_across\_populations} (default \texttt{.true.}). When set to \texttt{.false.}, the \texttt{get\_male\_from\_cell} function is called with \texttt{only\_population=jp}, enforcing strict reproductive isolation so that females can only mate with males of the same population.
\end{itemize}
\item \textbf{Files Modified}: \texttt{python\_interface.f95}, \texttt{mod\_agent\_world.f95}, \texttt{mod\_reviewed\_modules.f95}, \texttt{mod\_birth\_death\_new.f95}, \texttt{mod\_config.f95}, \texttt{mod\_read\_inputs.f95}, \texttt{basic\_config.nml}, \texttt{config\_sandesh.nml}, \texttt{simulation.py}, \texttt{application.py}.
\end{itemize}

\subsection{9. Critical Bugfix: Hashmap Population Corruption (May 26th)}
\begin{itemize}
\item \textbf{Objective}: Fix a bug causing phantom agents of incorrect populations (e.g. Pop 1 red dots) to appear in single-population simulations.
\item \textbf{Root Cause}: The \texttt{remove()} subroutine in \texttt{mod\_hashmap.f95} used \textbf{open-address linear probing}. When an entry was deleted, subsequent entries in the probe chain were rehashed to fill the hole. The rehash code correctly saved and restored the \texttt{key} and \texttt{value} (array index) fields, but \textbf{failed to save or restore the \texttt{population} field}. The moved entry inherited a stale \texttt{population} value from the vacated bucket.
\item \textbf{Impact}: When \texttt{get\_agent(id, world)} later looked up a rehashed agent via \texttt{get\_index\_and\_pop}, it received an incorrect population index. This caused an \textbf{out-of-bounds array access} into \texttt{agents(index, stale\_population)}, reading garbage memory that could appear as an agent with \texttt{population = 1}, rendered as a red dot in the visualization.
\item \textbf{Fix Applied} (4 changes to \texttt{mod\_hashmap.f95}):
\begin{itemize}
\item Added \texttt{temp\_population} variable to the rehash loop.
\item Clear \texttt{population = -1} when initially removing an entry's bucket.
\item Save \texttt{temp\_population} alongside \texttt{temp\_key}/\texttt{temp\_value} when reading entries to move.
\item Write \texttt{temp\_population} to the destination bucket and clear \texttt{population = -1} on the source bucket after moving.
\end{itemize}
\item \textbf{File Modified}: \texttt{src/data\_structures/mod\_hashmap.f95}.
\end{itemize}

\subsection{10. Efficient Cluster-Restricted Density Updates (May 26th)}
\begin{itemize}
\item \textbf{Objective}: Speed up the main bottleneck of the simulation step (the $O(N)$ grid sweeps for raw density, combined agent flows, box smoothing, and carrying capacity copies) by restricting updates to cells that belong to active clusters.
\item \textbf{Work Done}:
\begin{itemize}
\item \textbf{New Subroutine \texttt{update\_density\_and\_hep\_grid\_efficient}}: Implemented a cluster-restricted variant of the full density update in \texttt{mod\_technical\_modules.f95}. Instead of sweeping all $n_x \times n_y$ cells, it iterates only over cells belonging to active clusters (\texttt{cluster\_store\%clusters(k)\%cell\_gx/gy}), computing raw density, agent flows, and HEP copies exclusively for those cells.
\item \textbf{Hybrid Scheduling in \texttt{step\_simulation}}: Modified the density update dispatch in \texttt{python\_interface.f95} to choose between the full and efficient variants based on tick alignment. On clustering ticks (\texttt{mod(t, update\_interval) == 0}), the legacy \texttt{update\_density\_and\_hep\_grid} runs to ensure the full grid is populated before re-clustering. On intermediate ticks, the efficient variant is used.
\item \textbf{Boundary Orphan Reset}: Designed and implemented a \texttt{was\_clustered(:,:)} boundary logical grid array inside the \texttt{Grid} type. On re-clustering ticks (\texttt{mod(t, update\_interval) == 0}), any cell that transitioned OUT of a cluster is immediately zeroed out.
\item \textbf{O(1) Neighborhood box smoothing}: Restricted the smoothed density box-filter calculations to cluster cells, looking up neighbor cells directly (which naturally evaluate to \texttt{0.0} outside clusters).
\item \textbf{O(1) HEP data copies}: Restricted carrying capacity copies strictly to cells within active clusters.
\item \textbf{Conditional Startup Fallback}: Embedded conditional fallback inside \texttt{step\_simulation} in \texttt{python\_interface.f95} to automatically pivot to the legacy \texttt{update\_density\_and\_hep\_grid} on tick 0 or when no clusters exist, ensuring robust carrying capacity (\texttt{NC}) initialization and preventing agent die-offs.
\item \textbf{Config Flag}: Gated behind \texttt{efficient\_density\_updates} (boolean) in the namelist. Currently set to \texttt{.false.} in both \texttt{basic\_config.nml} and \texttt{config\_sandesh.nml} pending validation that it does not affect simulation accuracy.
\end{itemize}
\item \textbf{Status}: Implemented and compiles, but \textbf{currently disabled} (\texttt{efficient\_density\_updates = .false.}).
\item \textbf{Files Modified}: \texttt{python\_interface.f95}, \texttt{mod\_technical\_modules.f95}, \texttt{mod\_grid\_id.f95}, \texttt{mod\_config.f95}, \texttt{mod\_read\_inputs.f95}, \texttt{basic\_config.nml}, \texttt{config\_sandesh.nml}.
\end{itemize}

\subsection{11. Responsive Rolling Average for Performance HUD (May 26th)}
\begin{itemize}
\item \textbf{Objective}: Upgrade the performance metrics HUD from a global cumulative average to a responsive rolling average that forgets the distant past and immediately reacts to mid-run toggle switches.
\item \textbf{Work Done}:
\begin{itemize}
\item \textbf{20-Tick Sliding Window Buffer}: Designed and implemented a module-level sliding circular buffer array \texttt{perf\_history(6, 20)} in \texttt{python\_interface.f95} that records the exact timings of the last 20 ticks.
\item \textbf{Reset on Toggle}: Programmed the rolling timing history array to completely clear and restart fresh whenever performance timing is enabled or disabled in \texttt{set\_performance\_timing\_enabled}.
\item \textbf{Arithmetic Average stats}: Updated \texttt{get\_performance\_stats} to calculate the exact, unweighted arithmetic average of the recorded window (up to 20 samples), returning clean rolling averages without any clearing side-effects.
\end{itemize}
\item \textbf{Files Modified}: \texttt{src/interfaces/python\_interface.f95}.
\end{itemize}

\subsection{12. HEP Data Chunked Loading (May 26th)}
\begin{itemize}
\item \textbf{Objective}: Reduce memory footprint for large HEP NetCDF files ($\approx 900$~MB) by loading only a window of time slices instead of the entire dataset.
\item \textbf{Work Done}:
\begin{itemize}
\item \textbf{Dynamic Chunk Sizing}: \texttt{read\_hep\_data} in \texttt{mod\_read\_inputs.f95} now queries the HEP file size and calculates a \texttt{slices\_per\_chunk} count targeting $\approx 10$~MB chunks. Only the first chunk is loaded at startup.
\item \textbf{On-Demand Re-Loading}: Added \texttt{load\_hep\_chunk\_from\_file} subroutine. When \texttt{t\_hep} advances past the current chunk's range (\texttt{chunk\_start\_t..chunk\_end\_t}), a new chunk is loaded from the NetCDF file with the correct \texttt{start} offset.
\item \textbf{Local Index Translation}: \texttt{update\_density\_and\_hep\_grid} translates the global \texttt{t\_hep} to a local chunk index via \texttt{local\_idx = t\_hep - chunk\_start\_t + 1} before reading \texttt{grid\%hep(:,:,jp,local\_idx)}.
\item \textbf{Grid Type Extension}: Added \texttt{slices\_per\_chunk}, \texttt{chunk\_start\_t}, \texttt{chunk\_end\_t} metadata fields to the \texttt{Grid} type in \texttt{mod\_grid\_id.f95}, and a \texttt{config} pointer so the chunk loader can access file paths.
\end{itemize}
\item \textbf{Files Modified}: \texttt{mod\_read\_inputs.f95}, \texttt{mod\_grid\_id.f95}, \texttt{mod\_agent\_world.f95}, \texttt{mod\_technical\_modules.f95}.
\end{itemize}

\subsection{13. Clustering Fix: ``Clusters Way Too Small'' (May 26th)}
\begin{itemize}
\item \textbf{Objective}: Diagnose and fix the issue where watershed clusters were spatially too tight, capturing only $\approx 55\%$ of alive agents and missing $\approx 45\%$ in peripheral cells.
\item \textbf{Root Cause}: Three interacting factors:
\begin{enumerate}
\item \texttt{watershed\_threshold} was changed from \texttt{0.00} $\rightarrow$ \texttt{0.05}: With the old value, every non-zero-density cell was included. The new value excluded peripheral cells where smoothed density fell below \texttt{0.05}.
\item \texttt{human\_density\_smoothing\_radius} was changed from \texttt{3} $\rightarrow$ \texttt{5}: The wider kernel spreads density further, reducing peak values and causing more edge cells to drop below threshold.
\item Redundant internal smoothing in \texttt{watershed\_cluster()}: The watershed algorithm was applying 4 additional iterations of box-filter smoothing (radius=4) on top of the already-smoothed \texttt{human\_density\_smoothed} surface.
\end{enumerate}
\item \textbf{Fix Applied}:
\begin{itemize}
\item Lowered \texttt{watershed\_threshold} to \texttt{0.01} in both \texttt{basic\_config.nml} and \texttt{config\_sandesh.nml}.
\item Removed the redundant 4-iteration internal smoothing loop from \texttt{watershed\_cluster()} in \texttt{mod\_watershed.f95}.
\end{itemize}
\item \textbf{Files Modified}: \texttt{mod\_watershed.f95}, \texttt{basic\_config.nml}, \texttt{config\_sandesh.nml}.
\end{itemize}

\subsection{14. Bugfix: Out-of-Bounds Array Access in Cluster Accumulators (May 26th)}
\begin{itemize}
\item \textbf{Objective}: Fix memory corruption caused by accessing \texttt{cell\_cluster\_idx} with invalid grid coordinates.
\item \textbf{Root Cause}: Agents have a default \texttt{gx = -1} before being placed on the grid. The accumulator code in \texttt{update\_agent\_age} (\texttt{mod\_technical\_modules.f95}) accessed \texttt{cell\_cluster\_idx(agent\%gx, agent\%gy)} without bounds-checking. With \texttt{gx = -1}, this is an \textbf{out-of-bounds Fortran array access} that reads arbitrary memory.
\item \textbf{Same bug in debug code}: The detailed debug breakup prints in \texttt{python\_interface.f95} had the identical issue.
\item \textbf{Fix Applied}:
\begin{itemize}
\item Added \texttt{gx >= 1 .and. gx <= nx .and. gy >= 1 .and. gy <= ny} bounds check in \texttt{update\_agent\_age} before accessing \texttt{cell\_cluster\_idx}.
\item Added \texttt{c\_idx <= n\_clusters} upper-bound check before using \texttt{c\_idx} as an index into \texttt{clusters(:)}.
\item Added matching bounds checks in the debug breakup code in \texttt{python\_interface.f95}.
\end{itemize}
\item \textbf{Files Modified}: \texttt{mod\_technical\_modules.f95}, \texttt{python\_interface.f95}.
\end{itemize}

\section{Tripwires \& Lessons Learned}
\begin{itemize}
\item \textbf{1. Fortran Array Out-of-Bounds is Silent}:
Fortran does \textbf{not} bounds-check array accesses by default. Accessing \texttt{array(-1, 5)} doesn't crash --- it silently reads garbage from adjacent memory. This garbage then propagates through the simulation as valid data.
\item \textbf{2. Config Changes Can Have Non-Obvious Algorithmic Interactions}:
The \texttt{watershed\_threshold} increase from \texttt{0.00} to \texttt{0.05} seemed harmless --- but it interacted catastrophically with the independently-added 4x internal smoothing in the watershed algorithm.
\item \textbf{3. ``Parity'' Code Can Be Counter-Productive}:
The internal smoothing loop was added to bring watershed clustering into ``parity'' with Auto K-Means. But the two algorithms have fundamentally different sensitivities.
\item \textbf{4. Misleading Debug Output Can Send You in Circles}:
The debug diagnostic showing ``316 agents with invalid/unallocated indices'' looked like a clustering algorithm failure. In reality, it was the debug code itself having an out-of-bounds bug.
\item \textbf{5. \texttt{efficient\_density\_updates = .false.} Made Fix Attempts Into No-Ops}:
The earlier fix attempt (hybrid density scheduling, \texttt{was\_clustered} grid cleanup) was architecturally sound but was entirely guarded by \texttt{efficient\_density\_updates}, which was \texttt{.false.} in both configs.
\end{itemize}